
%...........................................................
%
% Comparison of Run 1-7b Asimov A with sensitivity
%
%...........................................................


\section {Comparison of Run 1-7b Asimov data set A with mean sensitivity}
\label{app:brazilian_asimov_mean_dcp_run17b}


In this section the mean sensitivity for $\delta_{CP} = -\pi /2$ and normal hierarchy
is compared with the Asimov data set A one ($\delta_{CP} = -1.601$ and normal hierarchy) 
for the Run 1-7b statistic.
The same procedure as described in \cref{sec:comparison_data_sensi_dcp_run17b} has been followed
by replacing the Run 1-7b data set with the Asimov data set A one.
In \cref{fig:brazilian_mpiover2_Asimov_A} it is shown that the mean sensitivity is very 
similar to the Asimov data set A one 
for values of $\delta_{CP}$ not far from the best-fit point,
while at $\delta_{CP} = +\pi/2$, for both hierarchies, the two expected 
$\Delta \chi^2$ distributions differ by about 1.
However it is worth to mention that the mean sensitivity is the result of the marginalization
of all the systematic parameters, oscillation included,
as well as of the statistical fluctuations,
while the Asimov data set A has 
all the oscillation parameters fixed to some arbitrary values
and the flux, cross section and detector systematic parameters fixed at their nominal values.
For this reason we do not expect mean and Asimov sensitivities to 
give the same results.

\begin{figure*}[htb]
\center
\begin{subfigure}[b]{0.49\textwidth}
\includegraphics[width=0.9\textwidth]{joint_plots/brazilian_plots/Run17b/brazilian_nh_dcp_mpiover2_Asimov_A.pdf}
\caption{Normal hierarchy}
\label{fig:brazilian_mpiover2_Asimov_A_nh}
\end{subfigure}
\begin{subfigure}[b]{0.49\textwidth}
\includegraphics[width=0.9\textwidth]{joint_plots/brazilian_plots/Run17b/brazilian_ih_dcp_mpiover2_Asimov_A.pdf}
\caption{Inverted hierarchy}
\label{fig:brazilian_mpiover2_Asimov_A_ih}
\end{subfigure}
\caption{
The comparison between the expected $<\Delta \chi^2>$ distribution 
for $\delta_{CP} = 0$ and normal hierarchy
and the Run 1-7b data set $\Delta \chi^2$ distribution is shown.
10k toy experiments are produced:
all the oscillation parameters which are not parameters of interest
are marginalised using the prior distributions shown in \cref{tab:osc_params_joint_priors}
and applying the reactor constraint.
The red solid line is the expected $<\Delta \chi^2>$ distribution.
The black solid line corresponds to the data $\Delta \chi^2$ distribution.
The green and light blue bands correspond to the regions
that contain respectively 68.27\% and 95.45\% of toy experiments
assuming the $\Delta \chi^2$ p.d.f. to be 1-sided.
The critical values, obtained with the Feldman-Cousins method,
that define the 90\% (blue) and $2\sigma$ (orange) confidence intervals
are shown as well.
}
\label{fig:brazilian_mpiover2_Asimov_A}
\end{figure*}



\clearpage



